% Part: first-order-logic
% Chapter: lindstrom
% Section: introduction

\documentclass[../../../include/open-logic-section]{subfiles}

\begin{document}

\olfileid[pt]{mod}{lin}{int}
\section{Introdução}

Neste capítulo, pretendemos demonstrar a caracterização de Lindstr\"om da
lógica de primeira ordem como a lógica máxima para a qual (dadas certas
restrições adicionais) valem os teoremas da Compacidade e de
L\"owenheim--Skolem Descendente
(\olref[fol][com][com]{thm:compactness} e
\olref[fol][com][dls]{thm:downward-ls}). Primeiro, precisamos de uma
caracterização mais geral da classe de lógicas à qual o teorema se aplica.
Vamos nos restringir a linguagens \emph{relacionais}, O.s.\@, linguagens que
contêm apenas !!{predicate}s e constantes individuais, mas nenhum
!!{function}.

\end{document}
